\section*{Chapter \ref{ch:b2:welfare-theorems} --- The Two Welfare Theorems}

\begin{solution}{exo:b2:welfare-theorems:1}
% code: exercises/solutions.py crusoe
The planner's condition $1/(2L) = 1/(16 - L)$ gives $L = 16/3 = 5.333$ and $c = 4\sqrt{5.333} = 9.238$.
The firm sets $w = 2/\sqrt L = 0.866$ and earns $9.238 - 0.866 \times 5.333 = 4.619$. The household's
full income is $16 \times 0.866 + 4.619 = 18.475$; it spends half on leisure, $18.475/(2 \times 0.866) =
10.667$ hours, so it works $16 - 10.667 = 5.333$ hours and consumes $9.238$.
\end{solution}

\begin{solution}{exo:b2:welfare-theorems:2}
% code: exercises/solutions.py ex_gap
% ledger: b2:persson2002-t3
The gap fell from 22.0 to 13.9 per cent, by 8.1 points; freight fell from 9.9 to 7.6, by 2.3 points, which
is $2.3/8.1 = 28$ per cent of the fall. The residual fell from 8.4 to 3.5 per cent: it more than halved.
\end{solution}

\begin{solution}{exo:b2:welfare-theorems:3}
% code: exercises/solutions.py ex_profit
Profit $3 \times 2\sqrt L - L$ is maximised where $3/\sqrt L = 1$: $L = 9$, $y = 6$, profit $18 - 9 = 9$.
\end{solution}

\begin{solution}{exo:b2:welfare-theorems:4}
Prices of the same quality in two markets between which the good is being shipped, with a gap that stays
above the full cost of shipping (freight, insurance, port charges, interest on the cargo during the
voyage) for long periods. Before concluding, one must check that the qualities are the same, that the
prices are for the same dates allowing for the voyage, that all costs are measured, and that trade was
actually flowing: without shipments the proposition allows any gap up to $t$.
\end{solution}

\begin{solution}{exo:b2:welfare-theorems:5}
% code: exercises/solutions.py crusoe
The planner maximises $a\ln(A\sqrt L) + (1 - a)\ln(T - L)$: $a/(2L) = (1 - a)/(T - L)$, so $a(T - L) = 2(1
- a)L$ and $L = aT/(2 - a)$. With $a = \tfrac23$: $L = (32/3)/(4/3) = 8$, $c = 4\sqrt8 = 11.31$ and $w =
2/\sqrt8 = 0.707$.
\end{solution}

\begin{solution}{exo:b2:welfare-theorems:6}
% code: exercises/solutions.py ex_tech
Hours do not depend on $A$ ($L = aT/(2 - a)$), so they stay at 5.333. Consumption and the real wage rise in
proportion to $A$, by 25 per cent: $c = 5\sqrt{5.333} = 11.547$ and $w = 5/(2\sqrt{5.333}) = 1.083$. With
log utility the higher wage makes leisure dearer (substitution effect, more work) and makes the household
richer (income effect, less work), and the two cancel.
\end{solution}

\begin{solution}{exo:b2:welfare-theorems:7}
(a) An external effect: the smoke enters the laundry's technology without a price. (b) Price taking: the
owner of the only source sets the price. (c) \omterm{def:b2:welfare-theorems:complete-markets}{Complete markets}: there is no market for grain delivered
in the state ``poor harvest''. (d) Convexity: a large fixed cost and a near-zero marginal cost make the
technology non-convex, and a price equal to marginal cost would not cover the cost of building.
\end{solution}

\begin{solution}{exo:b2:welfare-theorems:8}
% code: exercises/solutions.py ex_transfer
On the \omterm{def:b2:exchange-economies:contract-curve}{contract curve} $x_{A2} = 20 \times 6/16 = 7.5$, so $x_A = (6, 7.5)$ and $x_B = (4, 2.5)$. The
supporting price is $\mathit{MRS}_A = 7.5/6 = 1.25$. $A$'s bundle costs $1.25 \times 6 + 7.5 = 15$ and her
\omterm{def:b2:exchange-economies:endowment}{endowment} is worth $1.25 \times 8 + 2 = 12$: $B$ must give her 3 units of good 2. $B$'s wealth is then
$1.25 \times 2 + 8 - 3 = 7.5$, and he demands $x_{B1} = \tfrac23 \times 7.5/1.25 = 4$ and $x_{B2} =
\tfrac13 \times 7.5 = 2.5$, as required.
\end{solution}

\begin{solution}{exo:b2:welfare-theorems:9}
Farmers who would pay more than the expected cost of the indemnity plus the cost of running the insurance,
and potential insurers who could pool harvests whose failures are not all correlated: for instance,
regions whose harvests fail in different years, and farmers who cut consumption sharply after bad
harvests or hold costly precautionary stocks. If instead the bad harvests hit every farmer in every region
at once, or insurers cannot observe the harvest, the absence of the market may reflect a risk that cannot
be pooled or an information problem (\cref{ch:b2:adverse-selection}), not an inefficiency that a market
would remove.
\end{solution}

\begin{solution}{exo:b2:welfare-theorems:10}
Suppose $(x', y')$ is feasible, $\sum_i x'_i = \sum_i \omega_i + \sum_j y'_j$, every consumer weakly
prefers $x'_i$ and some consumer $k$ strictly. As in the exchange case, $p \cdot x'_k > p \cdot \omega_k +
\sum_j \theta_{kj}\,p \cdot y^*_j$ and $p \cdot x'_i \ge p \cdot \omega_i + \sum_j \theta_{ij}\,p \cdot
y^*_j$ for every $i$, by \omterm{def:b2:welfare-theorems:lns}{local non-satiation}. Summing and using $\sum_i \theta_{ij} = 1$: $p \cdot \sum_i
x'_i > p \cdot \sum_i \omega_i + \sum_j p \cdot y^*_j$. Each firm maximises profit at $p$, so $p \cdot y^*_j
\ge p \cdot y'_j$ for every $y'_j \in Y_j$. Hence $p \cdot \sum_i x'_i > p \cdot \sum_i \omega_i + \sum_j
p \cdot y'_j = p \cdot \sum_i x'_i$, a contradiction.
\end{solution}

\begin{solution}{exo:b2:welfare-theorems:11}
At $p = (1, 1)$ each consumer's wealth is 6. $A$ attains her maximum utility 4 with any affordable bundle
whose goods sum to at least 4, so $(3, 3)$ is a best choice; $B$ is indifferent among all bundles costing
6, so $(3, 3)$ is a best choice for him too. Markets clear with no trade: an equilibrium. Moving one unit of
good 1 from $A$ to $B$ gives $A$ $(2, 3)$, still worth 4 to her, and $B$ $(4, 3)$, worth $7 > 6$: a Pareto
improvement. The failing step is ``$p \cdot x'_i \ge p \cdot \omega_i$ for every consumer who is not
worse off'': $A$'s new bundle costs $5 < 6$ and is no worse for her, because she is satiated and \omterm{def:b2:welfare-theorems:lns}{local
non-satiation} fails.
\end{solution}

\begin{solution}{exo:b2:welfare-theorems:12}
% code: exercises/solutions.py ex_second
The target is on the \omterm{def:b2:exchange-economies:contract-curve}{contract curve}, so $x_{A2}$ follows from \eqref{eq:b2:exchange-economies:contract};
the supporting price is $A$'s marginal rate of substitution, $p = \frac{a}{1 - a}\,x_{A2}/x_{A1}$. $A$ spends
the share $a$ of her wealth $m_A$ on good 1, so she buys $x_{A1}$ exactly when $m_A = px_{A1}/a$. The
transfer to $A$ is $\tau_A = px_{A1}/a - (p\omega_{A1} + \omega_{A2})$, and $B$ gives it. With the numbers:
$x_{A2} = 12.5 \times 40/(0.375 \times 60 + 0.125 \times 40) = 500/27.5 = 18.18$, $p = 18.18/40 = 0.4545$,
$m_A = 0.4545 \times 40/0.5 = 36.36$, and $\tau_A = 36.36 - (0.4545 \times 80 + 10) = 36.36 - 46.36 = -10$:
$A$ gives 10 units of good 2.
\end{solution}

\begin{solution}{pb:b2:welfare-theorems:1}
% code: exercises/solutions.py problem
\textbf{Part I.}
1. If grain is shipped, $p_B = p_A + t$; if not, $p_B - p_A \le t$ (\cref{prop:b2:welfare-theorems:one-price}).
2. Ardel: $10p_A = 100 - 10p_A + s$. Brenmoor: $s = 200 - 10p_B$.
3. With $p_B = p_A + 6$: $s = 140 - 10p_A$, so $20p_A - 100 = 140 - 10p_A$ and $p_A = 8$, $p_B = 14$, $s =
60$ tonnes.
4. Without the road, $10p_A = 100 - 10p_A$ gives $p_A = 5$; Brenmoor consumes no grain.
5. $6 \times 60 = 360$ dalens a week.
6. Yes: it is a \omterm{def:b2:exchange-economies:competitive-equilibrium}{competitive equilibrium} with price-taking farms, households and carriers, no external
effects and a market for grain in each town, so the first welfare theorem applies. The freight is a real
cost of moving grain, not a waste the theorem ignores.
\textbf{Part II.}
7. Profit $p_A\sqrt{20L} - L$ is maximised where $p_A\sqrt{20}/(2\sqrt L) = 1$, so $L = 5p_A^2$ and $q =
\sqrt{100p_A^2} = 10p_A$.
8. $q = 80$ tonnes, $L = 5 \times 64 = 320$ hours, profit $8 \times 80 - 320 = 320$ dalens.
9. Ardel's households own the farms; the profit is part of their income, alongside their wages, and it is
what lets \omterm{def:b2:exchange-economies:walras-law}{Walras' law} hold in an economy with firms.
10. Ardel's consumers: $\tfrac12 \times 20 \times (10 - 8) = 20$. Brenmoor's: $\tfrac12 \times 60 \times (20
- 14) = 180$. Profit 320. Total 520 dalens.
\textbf{Part III.}
11. $s = 180 - 10p_A$ and $20p_A - 100 = 180 - 10p_A$: $p_A = 28/3 = 9.33$ and $p_B = 11.33$.
12. $s = 200 - 113.33 = 86.67$ tonnes; output $93.33$; Ardel consumes $100 - 93.33 = 6.67$.
13. $11.33 - 9.33 = 2 = t$.
14. In fractions, $p_A = 28/3$, $p_B = 34/3$ and $s = 260/3$. Ardel's consumers: $\tfrac12 \times \tfrac{20}{3}
\times \tfrac23 = 2.22$. Brenmoor's: $\tfrac12 \times \tfrac{260}{3} \times \tfrac{26}{3} = 375.56$. Profit:
$5 \times (28/3)^2 = 435.56$. Total $813.33$, a gain of $293.33$ dalens a week.
15. Ardel's consumers lose $20 - 2.22 = 17.78$: they now pay a price close to the city's.
16. Yes: the farms gain $115.56$ and Brenmoor's consumers $195.56$, far more than $17.78$. A \omterm{def:b2:welfare-theorems:lump-sum}{lump-sum
transfer}, fixed in advance and independent of how much grain anyone buys or grows, could compensate
Ardel's consumers without changing any choice.
17. In this model a gap above the transport cost is impossible when grain is shipped. A positive residual
would be either a cost the measure leaves out or a profit that competition has not removed.
18. A gap between the two prices close to the rail cost, prices that move together week by week, and
shipments that respond quickly when the gap widens.
\textbf{Part IV.}
19. $p = (\tfrac12 \times 10 + \tfrac14 \times 90)/(\tfrac12 \times 80 + \tfrac34 \times 20) = 27.5/55 = 0.5$
bolts per tonne. Ardel's wealth is $0.5 \times 80 + 10 = 50$: it consumes $(50, 25)$; Brenmoor $(50, 75)$.
20. Ardel's grain is 40, so its cloth is $12.5 \times 40/(0.375 \times 60 + 0.125 \times 40) = 18.18$;
Brenmoor gets $(60, 81.82)$.
21. $p = 18.18/40 = 0.4545$.
22. Ardel's bundle costs $0.4545 \times 40 + 18.18 = 36.36$; its \omterm{def:b2:exchange-economies:endowment}{endowment} is worth $0.4545 \times 80 + 10 =
46.36$. Ardel gives Brenmoor 10 bolts of cloth.
23. \omterm{def:b2:exchange-economies:endowment}{Endowments} become $(80, 0)$ and $(20, 100)$: $p = (\tfrac12 \times 0 + \tfrac14 \times 100)/55 = 0.4545$;
Ardel's wealth $36.36$ buys $\tfrac12 \times 36.36/0.4545 = 40$ tonnes and $18.18$ bolts. The target.
24. After the railway grain costs 9.33 dalens a tonne in Ardel and 11.33 in Brenmoor; a transfer of 10 bolts
of cloth from Ardel to Brenmoor delivers the council's \omterm{def:b2:exchange-economies:allocation}{allocation}.
\end{solution}
